O algoritmo Merge Sort apresenta comportamento semelhante nos casos
crescente, decrescente e aleatório, pois divide recursivamente o vetor
em duas partes e realiza a intercalação dos elementos em cada nível
da recursão.

\subsubsection{Melhor caso -- $\Theta(n\log n)$}

O melhor caso ocorre quando o vetor já está ordenado de forma crescente.
Mesmo nessa situação, o algoritmo continua dividindo o vetor em duas
metades até atingir intervalos com um único elemento.

Em seguida, as metades são intercaladas. Como a implementação realiza
a cópia dos elementos para vetores auxiliares e percorre as partes
durante a intercalação, o trabalho realizado em cada nível é
proporcional a $n$.

A quantidade de níveis de divisão é proporcional a $\log_2 n$.
Portanto, o custo total pode ser representado por:

\begin{equation}
T(n) = 2T\left(\frac{n}{2}\right) + cn
\end{equation}

em que $c$ representa uma constante associada ao trabalho de
intercalação.

Como existem aproximadamente $\log_2 n$ níveis, cada um com custo
proporcional a $n$, temos:

\begin{equation}
T(n) = \Theta(n\log n)
\end{equation}

Assim, o melhor caso do Merge Sort apresenta complexidade
linearítmica.

\subsubsection{Pior caso -- $\Theta(n\log n)$}

O pior caso ocorre quando a intercalação precisa comparar os elementos
das duas metades até que uma delas seja completamente percorrida.
Isso pode acontecer quando os elementos estão distribuídos de maneira
que seja necessário intercalar valores de ambas as metades.

Apesar do aumento do trabalho de comparação em cada intercalação, o
algoritmo continua dividindo o vetor em duas partes e processando todos
os elementos em cada nível.

Dessa forma, a recorrência permanece:

\begin{equation}
T(n) = 2T\left(\frac{n}{2}\right) + cn
\end{equation}

Como o custo por nível é proporcional a $n$ e existem
aproximadamente $\log_2 n$ níveis, temos:

\begin{equation}
T(n) = \Theta(n\log n)
\end{equation}

Portanto, o pior caso do Merge Sort também apresenta complexidade
linearítmica.

\subsubsection{Caso médio -- $\Theta(n\log n)$}

No caso médio, considera-se uma distribuição aleatória dos elementos
do vetor. O algoritmo realiza as divisões recursivas da mesma forma
e intercala as partes ordenadas até completar a ordenação.

Independentemente da distribuição inicial, todos os elementos precisam
ser processados durante as intercalações de cada nível. Assim, o
trabalho total continua sendo proporcional a $n$ por nível.

Consequentemente, temos:

\begin{equation}
T(n) = \Theta(n\log n)
\end{equation}

Portanto, o caso médio do Merge Sort apresenta complexidade
linearítmica.

\subsubsection{Resumo das complexidades}

\begin{table}[H]
    \centering
    \caption{Complexidades do algoritmo Merge Sort}
    \label{tab:complexidade-merge}
    \begin{tabular}{lc}
        \hline
        \textbf{Caso} & \textbf{Complexidade de tempo} \\
        \hline
        Melhor caso & $\Theta(n\log n)$ \\
        Caso médio & $\Theta(n\log n)$ \\
        Pior caso & $\Theta(n\log n)$ \\
        \hline
    \end{tabular}
\end{table}

\quad Conclui-se que o Merge Sort mantém complexidade temporal
$\Theta(n\log n)$ nos três casos analisados. Entretanto, sua
implementação convencional necessita de memória auxiliar proporcional
ao tamanho do vetor.
